\documentclass[12pt,aspectratio=169]{beamer}
\input{header_beamer}
\usetheme{Boadilla}
\usecolortheme{beaver}
\setbeamertemplate{page number in head/foot}{}
\setbeamertemplate{navigation symbols}{}
\title{Lecture-18: Stopping Times}
\date{}

\begin{document}
\pagestyle{empty}
\maketitle



\begin{frame}{Stopping times}
Let $(\Omega, \sF, P)$ be a probability space. %, and $T\subseteq \R$ an ordered index set considered as time. 
Consider a random process $X:\Omega\to\sX^T$ defined on this probability space with state space $\sX \subseteq \R$ and ordered index set $T\subseteq \R$ considered as time. 
\begin{Definition}<2->
A collection of event spaces denoted $\sG_{\bullet} \triangleq (\sG_t \subseteq \sF: t \in T)$ is called a filtration if $\sG_s \subseteq \sG_t$ for all $s \le t$. 
\end{Definition}
\begin{block}{Reminder}<3->
For the random process $X:\Omega\to\sX^T$, 
we can find the event space generated by all random variables until time $t$ as $\sF_t\triangleq \sigma(X_s, s \le t)$.  
The collection of event spaces $\sF_\bullet \triangleq (\sF_t: t \in T)$ is a filtration.
\end{block}
\end{frame}

\begin{frame}
\begin{Definition}
The natural filtration associated with a random process $X:\Omega \to \sX^T$ is given by $\sF_\bullet \triangleq (\sF_t: t\in T)$ where $\sF_t \triangleq \sigma(X_s, s\le t)$.  
\end{Definition}
\begin{block}{Reminder}<2->
For a random sequence $X: \Omega \to \sX^\N$, the natural filtration is a sequence $\sF_\bullet = (\sF_n: n \in \N)$ of event spaces $\sF_n \triangleq \sigma(X_1, \dots, X_n)$ for all $n \in \N$. 
\end{block}
\begin{block}{Reminder}<3->
If the random sequence $X$ is independent, 
then the random sequence $(X_{n+j}: j \in \N)$ is independent of the event space $\sigma(X_1, \dots, X_n)$. 
\end{block}
\end{frame}

\begin{frame}
\begin{Example} 
For a random walk $S$ with step size sequence $X$, 
the natural filtration of the random walk is identical to that of the step size sequence. 
That is, $\sigma(S_1, \dots, S_n) = \sigma(X_1, \dots, X_n)$ for all $n \in \N$. 
This follows from the fact that for all $n \in \N$, 
we can can write $S_j = \sum_{i=1}^jX_i$ and $X_j = S_j - S_{j-1}$ for all $j \in [n]$. 
That is, there is a bijection between $(X_1, \dots, X_n)$ and $(S_1, \dots, S_n)$. 
\end{Example}
\end{frame}

\begin{frame}
\begin{Definition}
\label{defn:StoppingTimeGen}
A random variable $\tau: \Omega \to T$ is called a \textbf{stopping time} with respect to  a filtration $\sF_\bullet$ if 
\begin{enumerate}[(a)]
\item the event $\tau^{-1}(-\infty, t] \in \sF_t$ for all $t \in T$, and 
\item the random variable $\tau$ is finite almost surely, i.e. $P\set{\tau < \infty} = 1$.  
\end{enumerate}
\end{Definition}
\end{frame}

\begin{frame}
\begin{block}{Reminder}
Intuitively, if we observe the process $X$ sequentially, 
then the event $\set{\tau \le t}$ can be completely determined by the observation $(X_s, s \le t)$ until time $t$. 
The intuition behind a stopping time is that its realization is determined by the past and present events but not by future events. 
That is, given the history of the process until time $t$, we can tell whether the stopping time is less than or  equal to $t$ or not. 
In particular, $\E[\SetIn{\tau \le t}\given \sF_t]=\SetIn{\tau \le t}$ is either one or zero. 
\end{block}
\end{frame}

\begin{frame}
\begin{Definition}[First hitting time]
For a process $X:\Omega\to\sX^T$ and any Borel measurable set $A \in \sB(\sX)$, we define the \textbf{first hitting time} $\tau_X^A: \Omega \to T\cup\set{\infty}$ for the process $X$ to hit states in $A$, as 
$\tau_X^A \triangleq \inf\set{t \in T: X_t \in A}.$
\end{Definition}

\begin{Example}<2->
We observe that the event $\set{\tau_X^A \le t} = \set{X_s \in A \text{ for some }s \le t} \in\sF_t$ for all $t \in T$. 
It follows that, if $\tau_A$ is finite almost surely, then $\tau_A$ is a stopping time with respect to filtration $\sF_\bullet$.
\end{Example}
\end{frame}


\begin{frame}
\begin{prop} 
For a random sequence $X: \Omega \to \sX^\N$, 
an almost sure finite discrete random variable $\tau: \Omega \to \N\cup\set{\infty}$ is a \textbf{stopping time} with respect to this random sequence $X$ iff the event $\set{\tau = n} \in \sigma(X_1, \dots, X_n)$ for all $n \in \N$. 
\end{prop}
\begin{Proof}<2->
From Definition~\ref{defn:StoppingTimeGen}, 
we have $\set{\tau = n} = \set{\tau \le n}\setminus\set{\tau \le n-1} \in \sigma(X_1, \dots, X_n)$. 
Conversely, from the theorem hypothesis, it follows that $\set{\tau \le n} = \cup_{m=1}^n\set{\tau = m} \in \sigma(X_1, \dots, X_n)$. 
\end{Proof}
\end{frame}


\begin{frame}
\begin{Example} 
Consider a random sequence $X: \Omega \to \sX^\N$, 
with the natural filtration $\sF_\bullet$,   
and a measurable set $A \in \sB(\sX)$. 
If the first hitting time $\tau_X^A:\Omega\to\N\cup\set{\infty}$ for the sequence $X$ to hit set $A$ is 
almost surely finite, then $\tau_X^A$ is a stopping time. 
This follows from the fact that $\set{\tau_X^A = n} = \cap_{k=1}^{n-1}\set{X_k \notin A}\cap\set{X_n \in A} \in \sF_n$ for each $n \in \N$. 
\end{Example}
\end{frame}

\begin{frame}
\begin{Definition} 
Consider a random process $X:\Omega\to \sX^{\R_+}$ with discrete state space $\sX\subseteq\R$. 
For each state $y \in \sX$, we define $\tau_X^{\set{y},0} \triangleq 0$ and inductively define the \textbf{$k$th hitting time to a state $y$} after time $t = 0$, as 
\EQ{
\tau_X^{\set{y},k} \triangleq \inf\set{t > \tau_X^{\set{y}, k-1}: X_t = y},\quad k \in \N. 
} 
\end{Definition}
\begin{block}{Reminder}<2->
We observe that $\set{\tau_X^{\set{y}, k} \le t} \in \sF_t$ for all times $t \in \R_+$. 
Hence if $\tau_X^{\set{y}, k}$ is almost surely finite, then it is a stopping time for the process $X$. 
\end{block}
\end{frame}


\begin{frame}
\begin{Definition} 
For a discrete valued random sequence $X:\Omega\to\sX^\N$, the number of visits to a state $y \in \sX$ in first $n$ time steps is defined as $N_y(n) \triangleq \sum_{k=1}^n\SetIn{X_k = y}$ for all $n \in \N$. 
\end{Definition}
\begin{block}{Reminder}<2->
We observe that $N_y:\Omega\to\Z_+^{\Z_+}$ is a random walk with the Bernoulli step size sequence $(\SetIn{X_k = y}: k \in \N)$. 
Further, $\tau_X^{\set{y}, k} = \tau_{N_y}^{\set{k}} = \inf\set{n \in \N: N_y(n) = k}$. 
We also observe that $\{N_y(n) \le k\} = \{\tau_y^{(k+1)} > n\}$ and $\{N_y(n) = k\} = \{\tau_y^{k} \le n < \tau_y^{(k+1)}\}$. 
\end{block}
\end{frame}

\begin{frame}
\begin{block}{Reminder}
We observe that the number of visits to state $y$ in first $n$ steps of $X$ is also given by 
\EQ{
N_y(n) 
= \sup\set{k \in \Z_+: \tau_X^{\set{y},k} \le n}
= \inf\set{k \in \N: \tau_X^{\set{y},k} > n} - 1
= \sum_{k\in\N}\SetIn{\tau_X^{\set{y},k} \le n}.
}
This implies that $N_y(n)+1$ is the first hitting time to set of states $\set{n+1, n+2, \dots }$ for the increasing random sequence $(\tau_X^{\set{y},k}:k \in \N)$. 
\end{block}
\end{frame}


\begin{frame}
\begin{block}{Lemma : Wald's Lemma} 
Consider a random walk $S: \Omega \to \R^{\Z_+}$ with \iid step-sizes $X: \Omega \to  \R^\N$ having finite $\E\abs{X_1}$. 
Let $\tau$ be a finite mean stopping time with respect to this random walk. 
Then,
\EQ{
\expect{S_\tau}= \expect{X_1}\expect{\tau}.
}
\end{block}
\end{frame}

\begin{frame}
\begin{Proof} 
Recall that the event space generate by the random walk and the step-sizes are identical. 
From the independence of step sizes, it follows that $X_n$ is independent of $\sigma(X_0, X_1, \dots, X_{n-1})$. 
Since $\tau$ is a stopping time with respect to random walk $S$, 
we observe that $\set{\tau \ge n} = \set{\tau > n-1} \in \sigma(X_0, X_1, \dots, X_{n-1})$,  
and hence it follows that random variable $X_n$ and indicator $\SetIn{\tau \ge n}$ are independent and $\E[X_n\SetIn{\tau \ge n}] = \E X_1 \E\SetIn{\tau \ge n}$. 
Therefore,
\begin{align*}
\E[S_\tau]
&= \E[\sum_{n=1}^\tau X_n]
= \E[\sum_{n \in \N} X_n\SetIn{\tau \ge n}]
= \sum_{n \in \N} \E X_1 \E\left[\SetIn{\tau \ge n}\right] \\
&= \E X_1\E\left[ \sum_{n \in \N}\SetIn{\tau \ge n}\right] 
= \E[X_1]\E[\tau].
\end{align*}
We exchanged limit and expectation in the above step, which is not always allowed. 
We were able to do it %since the summand is positive and we apply monotone convergence theorem. 
by the application of dominated convergence theorem. 
\end{Proof}
\end{frame}


\begin{frame}
\begin{cor} 
Consider the stopping time $\tau_S^{\set{i}} \triangleq \min\set{n \in \N: S_n = i}$ for an integer random walk $S: \Omega \to \Z_+^{\Z_+}$ with \iid step size sequence $X: \Omega \to \Z^\N$.  
Then, the mean of stopping time $\E \tau_S^{\set{i}} = {i}/{\E X_1}$. 
\end{cor}
\begin{Proof}<2->
This follows from the Wald's Lemma and the fact that $S_{\tau_i} = i$. 
\end{Proof}
\end{frame}


\begin{frame}{Properties of stopping time} 
\begin{block}{Lemma}
Let $\tau_1,\tau_2$ be two stopping times with respect to filtration $\sF_\bullet$.  
Then the following hold true. 
\begin{enumerate}[i\_]
\item $\min \set{\tau_1,\tau_2} $ is a stopping time. 
\item If $T$ is separable, then $\tau_1+\tau_2$ is a stopping time.
\end{enumerate}
\end{block}
\end{frame}


\begin{frame}
\begin{Proof}
Let $\tau_1,\tau_2$ be stopping times with respect to a filtration $\sF_{\bullet} = (\sF_t : t \in T)$.  
\begin{enumerate}[i\_]
\item<1-> Result follows since the event $\set{ \min \set{ \tau_1,\tau_2} > t} = \set{\tau_1 > t} \cap \set{\tau_2 > t} \in \sF_t$. 
\item<2-> A topological space is called separable if it contains a countable dense set. 
Since $\R_+$ is separable and ordered, we assume $T = \R_+$ without any loss of generality. 
It suffices to show that the event $\set{\tau_1+\tau_2 \le  t} \in \sF_t$ for $T = \R_+$. 
To this end, we observe that 
\EQ{
\set{\tau_1 + \tau_2 \le t} = \bigcup_{s \in \Q_+:~ s \le t}\set{\tau_1 \le t - s, \tau_2 \le s} \in \sF_t.
}
\end{enumerate}
\end{Proof}
\end{frame}


\begin{frame}{Strong independence property and applications}
\begin{block}{Theorem : Strong independence property} Let $X: \Omega\to\sX^{\R_+}$ be an independent random process with natural filtration $\sF_\bullet$, and $\tau:\Omega\to\R_+$ a stopping time with respect to $\sF_\bullet$, 
then $(X_{\tau+s}: s \in \R_+)$ is independent of history $(X_s: s \le \tau)$. 
\end{block}

\begin{Definition}<2-> 
We can define the \textbf{$k$th return time to state $y$} for the random process $X:\Omega\to\sX^{\R_+}$ as the interval between two successive visits to state $y$, that is for all $k \in \N$
\EQ{
H_X^{\set{y}, k} 
\triangleq \tau_X^{\set{y},k} - \tau_X^{\set{y},k-1}
= \inf\set{s \in \R_+: X_{\tau_X^{\set{y},k-1}+s} = y}.
}
\end{Definition}
\end{frame}


\begin{frame}
\begin{block}{Reminder}
We observe that $\tau_X^{\set{y}, k} = \sum_{j=1}^kH_X^{\set{y}, j}$ for all $k \in \N$. 
That is, the hitting times sequence $(\tau_X^{\set{y},k}: k \in \N)$ is a random walk with step-size sequence being return times $(H_X^{\set{y},j}: j \in \N)$. 
Therefore, if $H_X^{\set{y}, j}$ is almost sure finite for all $j \in [k]$, 
then the finite sum $\tau_X^{\set{y}, k}$ is almost sure finite for all $k\in \N$.  
\end{block}
\begin{block}{Reminder}<2-> 
From the bijection between hitting and return times, 
we observe that $\sigma(\tau_X^{\set{y}, k)}: k \in [n]) =  \sigma(H_X^{\set{y}, j}: j \in [n])$. 
Recall that $\tau_X^{N_y(n)+1} > n$ by definition. 
If passage times are independent and \iid from second passage time, 
then it follows from the Wald's Lemma that 
\EQ{
\E H_X^{\set{y},1} + \E[N_y(n)]\E H_X^{\set{y},2} \ge (n+1).
}
\end{block}
\end{frame}


\begin {frame}
\begin{Example}
We also observe that the $k$th hitting time to $\set{1}$ by a Bernoulli step size sequence $X:\Omega\to\set{0,1}^\N$ is the first hitting time to $\set{k}$ by random walk $S:\Omega\to\Z_+^{\Z_+}$. 
That is,  
\EQ{
\tau_X^{\set{1},k}
= \tau_S^{\set{k}} 
= \inf\set{n \in \N: S_n = k} 
= \tau_S^{\set{k-1}} + \inf\set{n \in \N: S_{\tau_S^{\set{k-1}}+n}-S_{\tau_S^{\set{k-1}}} = 1}. 
}
\end{Example}
\end{frame}


\begin{frame}
\begin{block}{Lemma} 
For an \iid Bernoulli random sequence $X:\Omega\to\set{0,1}^\N$ with $\E X_1 \in (0,1)$, 
the $k$th hitting time to state~$1$ is a stopping time, and $\tau_X^{\set{1}, k} = \sum_{i=1}^kY_i$ where $Y:\Omega\to\N^\N$ is an \iid random sequence distributed identically to $\tau_X^{\set{1}}$ . 
\end{block}

\begin{Proof}<2->
When the Bernoulli step size sequence $X$ is \iid with $\E X_1 = p \in (0,1)$, 
we get that $P\set{\tau_S^{\set{1}}= n} = (1-p)^{n-1}p$ for all $n \in \N$. 
It follows that 
\EQ{
P\set{\tau_S^{\set{1}}< \infty} 
= P\left(\cup_{n\in\N}\set{\tau_S^{\set{1}}=n}\right) 
= \sum_{n\in\N}P\set{\tau_S^{\set{1}}= n} 
= 1.
}
Hence, the random time $\tau_S^{\set{1}}$ is finite almost surely. 
\end{Proof}
\end{frame}


\begin{frame}
\begin{Proof}[Proof continued.]
We will show that $\tau_S^{\set{k}}$ is finite almost surely for all $k \in \N$ by induction. 
By induction hypothesis, $\tau_S^{\set{k-1}}$ is finite almost surely. 
Then $S_{\tau_S^{\set{k-1}}+n}-S_{\tau_S^{\set{k-1}}} = \sum_{j=1}^nX_{\tau_S^{\set{k-1}}+j}$ is the sum of $n$ \iid Bernoulli random variables independent of $\tau_S^{\set{k-1}}$ by strong independence property,  
and hence has distribution identical to $S_n$. 
Further,  
This implies that $\tau_S^{\set{k}} = \tau_S^{\set{k-1}} + \tau_S^{\set{1}}$, 
where $\tau_S^{\set{1}}$ has the identical distribution to $\tau_X^{\set{1}}$ and is independent of $\tau_S^{\set{k-1}}$. 
\end{Proof}   
\end{frame}

\end{document}
